\chapter{Active-Probing Battery Audit}
\label{app:active-probing-audit}

The main thesis now uses Chapter~\ref{ch:certos-runtime} for the runtime
promotion slot.  The active-probing material is therefore preserved here as
an evidence appendix rather than a mainline chapter.  That placement fits
the current claim boundary: the repo contains a bounded battery spoofing
detection package, but not a complete adversarial-safety frontier.

\section{Battery Threat Model}

The battery-specific adversary considered here attempts to preserve the
appearance of healthy telemetry while drifting away from true plant state.
Operationally, that means:

\begin{enumerate}
  \item replaying or smoothing readings so the battery appears stable,
  \item spoofing load-sensitive inputs so the dispatch policy overcommits,
    and
  \item hiding the attack inside packet behavior that otherwise appears
    plausible.
\end{enumerate}

The point of active probing is to break that illusion by asking the plant
for a small physical response that a forged telemetry stream cannot mimic
indefinitely.

\section{Locked Detection Package}

The bounded evidence package is:

\begin{itemize}
  \item \path{reports/probing/spoofing_attack_results.csv}, which stores the
    confusion counts, and
  \item \path{reports/probing/probing_detection_quality.png}, which turns the
    same audit into a defense-facing figure.
\end{itemize}

These outputs correspond to the publication-facing lineage rooted in
\path{reports/publication/active_probing_spoofing_detection.csv}.

\section{Detection Results}

\begin{table}[ht]
\centering
\caption{Locked active-probing detection results for the battery threat
model.}
\label{tab:active-probing-results}
\begin{tabular}{lr}
\toprule
\textbf{Metric} & \textbf{Value} \\
\midrule
Events evaluated & 300 \\
Threshold & 0.60 \\
True positives & 48 \\
False positives & 6 \\
False negatives & 4 \\
True negatives & 242 \\
Precision & 0.889 \\
Recall & 0.923 \\
\bottomrule
\end{tabular}
\end{table}

\section{Battery-Operational Meaning}

For the battery thesis, these numbers support a narrow operational point.
The probing layer catches most spoofing events while keeping false alarms
bounded.  That matters because every false alarm tightens the battery
unnecessarily, while every missed detection risks preserving the hidden
observation-action gap identified in Part~I.

\section{Bounded Interpretation}

This appendix supports three claims and no more:

\begin{itemize}
  \item the battery implementation contains an active-probing detection path,
  \item that path achieves high recall and strong precision on the current
    spoofing audit, and
  \item the audit does \emph{not} yet establish adversarial completeness or a
    full controller-vs.-attacker safety frontier.
\end{itemize}

Keeping this material in an appendix is therefore not a demotion of the
implementation.  It is a promotion of the manuscript's honesty.
